% Phase 1 -- written by Huzaifa
\section{Phase 1: Problem Definition}

\subsection{Problem Statement}
Tracking what we eat is one of the simplest ways to manage weight and health,
but it only works if logging a meal is quick. In South Asia it usually is not.
Popular nutrition apps are built around Western food databases, so everyday
dishes such as biryani, nihari, haleem and karahi are either missing, listed
under unfamiliar names, or represented by a generic entry that does not match
what is on the plate.

Users are left to search manually, guess which entry is closest, and estimate
the portion by eye. The error is large because the calorie content of desi
dishes depends heavily on the oil or ghee used, the cut of meat and the serving
size, none of which a generic entry captures. Most people give up after a few
days, which defeats the purpose of tracking.

Nutrify addresses this by turning a single photo of a meal into a dish name
and a calorie range. The system has three parts: an image classification model
that recognises 14 popular desi dishes, a calorie module that maps the
recognised dish to a per-serving value with an adjustable portion size, and a
web application with image upload and camera capture that returns the dish
name, a confidence score and a calorie range in seconds.

\subsection{Target Audience}
\textbf{Primary users.} Everyday people in Pakistan and the wider South Asian
region who want to keep an eye on what they eat, for example to lose or gain
weight, follow a fitness plan, or simply eat more mindfully. They are
comfortable taking a photo on a phone but do not want to type in meals or
look up nutrition tables.

\textbf{Secondary users.} Students, fitness coaches and dietitians who need a
quick estimate when discussing local food, and who currently rely on printed
tables or Western-centric apps.

\textbf{Research gap.} Widely used food recognition benchmarks such as
Food-101 are dominated by Western dishes, and there are few public, labelled
image collections for Pakistani food. As a result, models trained on those
benchmarks do not transfer well to dishes like nihari or haleem. This project
helps fill that gap by assembling a curated dataset of desi dishes, documenting
how it was collected and cleaned, and evaluating how the model behaves on
look-alike dishes, different lighting and restaurant versus home-cooked photos.

\subsection{AI Justification}
\textbf{Why not rule-based programming.} Recognising a dish from a photo cannot
be written as a set of if--else rules. The same dish looks different from one
cook, restaurant or camera to the next, and the lighting, plating and
background change from photo to photo. Dishes such as nihari, haleem and
aloo gosht are also visually similar brown stews, so simple colour or texture
rules would confuse them. Any hand-written rule set would need constant
exceptions and would still fail on photos it has not seen.

\textbf{Why deep learning fits.} Convolutional neural networks learn the
relevant visual features, from edges and textures up to whole-dish patterns,
directly from labelled examples instead of relying on features chosen by hand.
This is the standard approach for image classification and it handles the
variation in appearance that defeats fixed rules.

\textbf{Why transfer learning.} Collecting thousands of images per dish is not
realistic for a student project. Starting from a network pretrained on a large
general image collection (EfficientNet-B0) lets us reuse features it has
already learned and fine-tune it on roughly 150 to 300 images per dish. We will
compare it against a baseline CNN trained from scratch to show the benefit.
EfficientNet-B0 is also small enough to serve through a web application.

\textbf{Where AI stops.} Machine learning is used only for the part that needs
it, recognising the dish. The calorie value comes from a verified per-serving
entry in a food composition database, scaled by the chosen portion size. Keeping
that step deterministic makes the numbers traceable and easier to defend than a
model that guesses calories directly.
